\subsection{The global elliptic detectors necessarily have infinite ramification}

\begin{result}{E.200}{Proved}
{Conditional irreducibility and infinite ramification of the three cusp-specific detectors}
Let \((A,B)\) be an ordinary completed pair of infinite valuation
support.  Suppose that a full geometrically pro-three section over
\(\mathbb Q\) has its three-primary abelianization.  The three
pullbacks of the elliptic family E.192.2 described below give continuous
representations
\[
             \rho_c:G_{\mathbb Q}\longrightarrow
                         \operatorname{GL}_2(\mathbb Z_3),
             \qquad c\in\{0,1,\infty\}.
\tag{E.200.1}
\]
Each \(\rho_c\) is ramified at infinitely many finite primes.
At least two of these representations are absolutely irreducible.
More precisely, \(\rho_c\) is absolutely irreducible whenever the
ordinary support contains, away from three, both a prime in direction
\(c\) and a prime in one of the other two directions.

These are necessary conditions on a putative full section.  They do
not assert that such a section exists, and they do not exclude it.
In particular the detector does not turn the global problem into a
finite-ramification problem.
\end{result}

\begin{proof}
\emph{The three directions and their local elliptic types.}
For an ordinary pair put
\[
\begin{aligned}
 S_0&=\{q:v_q(A)>0,\ v_q(B)=0\},\\
 S_1&=\{q:v_q(A)=0,\ v_q(B)>0\},\\
 S_\infty&=\{q:v_q(A)=v_q(B)<0\}.
\end{aligned}
\tag{E.200.2}
\]
These are the only nonzero valuation patterns, since locally the
coordinates are \((t_q,1-t_q)\).  Their union is the valuation support.
Choose the tripod automorphisms
\[
 \phi_\infty(t)=t,\qquad
 \phi_0(t)=1/t,\qquad
 \phi_1(t)=1/(1-t).
\tag{E.200.3}
\]
Each sends its named cusp to infinity and the other two cusps to
zero and one.  Pull back E.192.2 along \(\phi_c\), and then along
the assumed section.  E.192 proves that its actual arithmetic
Tate-module representation factors through the geometrically
pro-three quotient; hence these are indeed E.200.1.  Their
determinants are \(\chi_3\).

We record the local reduction of E.192.2 over a finite extension of
\(\mathbb Q_p\), for \(p\ne3\).  At a parameter \(t\) with positive
valuation, the model has unit \(c_4\), discriminant valuation
\(3v(t)>0\), and the split node \(y(y+3x)=x^3\).  It is split
multiplicative.  At \(v(1-t)>0\), its \(c_4\) is again a unit and
the discriminant valuation is \(v(1-t)>0\).  On the special fibre,
put \(x=X-1\), \(y=Y+1\).  The equation becomes
\[
                    Y^2+3XY+3X^2=X^3.
\tag{E.200.4}
\]
Its two tangent slopes are the distinct roots of
\(r^2+3r+3\); they are \(\zeta_3-1\) and
\(\zeta_3^2-1\).  Thus this is multiplicative reduction, becoming
split over an unramified extension of degree at most two.
Both calculations remain valid at \(p=2\).

At \(v(t)<0\), adjoin \(s\) with \(s^3=t^{-1}\), and put
\(x=s^{-2}X\), \(y=s^{-3}Y\).  This gives the integral equation
\[
                       Y^2+3sXY+Y=X^3,
       \qquad \Delta=27(s^3-1).
\tag{E.200.5}
\]
Its discriminant is a unit, so the curve acquires good reduction
over this finite extension.  No restriction on \(v(t)\bmod3\)
is needed.  In particular the parameter-infinity direction is
potentially good, whereas the other two directions are
multiplicative.  Tate uniformization in the split cases is
\citet[C.14, Theorem 14.1, pp.444--445]{AuditSilverman2009}.

At every \(q\ne3\) in E.200.2, the full local section is the
ordinary point class by A.31.  The nonzero ordinary integer
valuation remains nonzero in \(\mathbb Z_3\), so that theorem's
ramified hypothesis applies.  It follows that the restriction of
\(\rho_c\) is the actual Tate-module representation of the elliptic
curve at \(\phi_c(t_q)\).  Therefore it has the local types just
calculated.  This invokes no point matching at zero-support places
or at the coefficient prime three.

\emph{The global character calculation.}
Let \(E/\mathbb Q_3\) be finite and let
\(\lambda:G_{\mathbb Q}\to E^\times\) be a continuous character.
Its compact image has valuation zero and lies in
\(\mathcal O_E^\times\).  Extend the usual three-adic logarithm
from principal units to this group, killing its finite torsion.
There is an \(a\in E\) such that
\[
                   \log\lambda(g)=a\log\chi_3(g)
                   \quad(g\in G_{\mathbb Q}).
\tag{E.200.6}
\]
Here is a proof which does not assume finite ramification.
Kronecker--Weber identifies
\(G_{\mathbb Q}^{\rm ab}=\widehat{\mathbb Z}^{\times}\)
\citep[I, Theorem 4.16, pp.49--50]{RootMilneCFT}.
For every prime \(r\ne3\), a continuous homomorphism from
\(\mathbb Z_r^\times\) to the additive group of \(E\) vanishes:
its finite torsion maps to zero, and a pro-\(r\) group has no
nontrivial compact pro-three image.  The finite-support product
of these factors is dense in their full product, so continuity
kills the entire product away from three.  On
\(\mathbb Z_3^\times\), logarithm identifies the principal units
with \(3\mathbb Z_3\).  A continuous additive homomorphism from
\(\mathbb Z_3\) to \(E\) is scalar multiplication, proving
E.200.6.  In particular, if \(a=i\in\{0,1\}\), then
\[
                       \lambda=\epsilon\chi_3^{\,i}
                       \quad\hbox{with \(\epsilon\) finite}.
\tag{E.200.7}
\]
Indeed the quotient has zero logarithm, and the kernel of logarithm
in \(\mathcal O_E^\times\) is its finite group of roots of unity.

\emph{Absolute irreducibility.}
Fix \(c\) and assume there are primes \(p\in S_c\setminus\{3\}\)
and \(q\in(S_0\cup S_1\cup S_\infty)\setminus(S_c\cup\{3\})\).
If \(\rho_c\) were reducible over \(\overline{\mathbb Q}_3\),
an invariant line would be defined over a finite extension \(E\),
and its character \(\lambda\) would satisfy E.200.6.  At \(q\),
pass to the unramified extension of residue degree \(f\le2\)
over which the elliptic curve is split multiplicative.  Every
arithmetic Frobenius lift has eigenvalues \(1,q^f\), by the Tate
exact sequence.  Thus
\[
                       a f\log q\ \in\ \{0,f\log q\}.
\]
Since the positive rational prime \(q\ne3\) is not a root of
unity in \(\mathbb Q_3^\times\), \(\log q\ne0\), and
\(a=0\) or \(1\).  Consequently E.200.7 holds.

At \(p\), pass instead to a finite extension over which
E.200.5 gives good reduction, and write its residue field order
as \(p^h\).  On an arithmetic Frobenius lift the character
E.200.7 has value a root of unity times \(p^{hi}\), with
\(i=0\) or \(1\).  On the other hand both eigenvalues of the
good elliptic fibre have, in every complex embedding, absolute
value \(p^{h/2}\)
\citep[V.2, Theorem 2.3.1, pp.142--143]{AuditSilverman2009}.
The value in question is algebraic over \(\mathbb Q\), since its
root-of-unity factor is algebraic.  Its complex absolute value
is \(1\) or \(p^h\), a contradiction.  This proves the stated
absolute irreducibility criterion.

\emph{Infinite ramification and the number of irreducible detectors.}
At every multiplicative support prime \(q\ne3\), the Tate
parameter \(Q\) after the unramified splitting extension has
positive valuation.  Its three-primary Kummer character on
inertia is
\(v_q(Q)\,t_3\), where \(t_3\) is the surjective three-primary
tame character.  This is nonzero in \(\mathbb Z_3\): a nonzero
ordinary integer cannot annihilate that module.  Thus \(\rho_c\)
is ramified at every prime of the other two support directions,
apart from the omitted prime three.  Divisibility of the
ordinary valuations by arbitrarily high powers of three at
successive primes does not change this conclusion.

At least two of the three sets E.200.2 must be infinite.  If only
\(S_0\) were infinite, \(B\) would have finite support; if only
\(S_1\) were infinite, \(A\) would have finite support; if only
\(S_\infty\) were infinite, \(A/B\) would have finite support.
Each alternative contradicts E.101, which makes such an
everywhere ordinary pair a rational pair of finite support.
It follows that the union of the other two directions is
infinite for every \(c\), proving infinite ramification of all
three representations.  For each of at least two infinite
directions \(c\), choose \(p\) there and \(q\) in another
infinite direction, avoiding three.  The irreducibility criterion
then applies to \(\rho_c\).

This argument supplies no contradiction to continuity of an
infinitely ramified three-adic representation.  Indeed the
irreducibility criterion can already hold for a rational point
with support in two directions.  A finite-ramification theorem
cannot be applied without an additional hypothesis establishing
its premise; the detector provides the opposite premise for
an ordinary infinite-support pair.  The ordinary support-density
conclusion of A.33 is independent of these representations.
\end{proof}
